\chapter{Vaginal dryness}
\index{Vaginal dryness}

\section*{Definition and Overview}
Vaginal dryness is a feeling of lack of moisture or lubrication in and around the vagina. It is caused by the vaginal lining becoming thin, less elastic, and less able to produce fluid, a state known medically as vaginal atrophy or atrophic vaginitis. It most commonly develops around and after menopause, when falling oestrogen levels change the vaginal tissues, but it can also occur during breastfeeding, after certain cancer treatments, and as a side effect of some medicines. Vaginal dryness can cause discomfort, soreness, and itching, and it can make sexual intercourse painful, which in turn can affect relationships and quality of life. It is a very common condition, yet many women do not raise it with their doctor because they feel embarrassed or assume it is an unavoidable part of ageing. In fact, there are many effective treatments, and no woman needs to simply put up with it.

\section*{Epidemiology}
Vaginal dryness affects a large proportion of women. Studies suggest that roughly half of postmenopausal women experience vaginal dryness at some point, making it one of the most common symptoms of menopause after hot flushes. Many women experience symptoms for years, but only a minority seek treatment, in part because of embarrassment. Vaginal dryness is also common in women who are breastfeeding, in whom the hormone changes that support milk production suppress oestrogen. It can occur in younger women who take certain hormonal contraceptives or anti-oestrogen medicines, and in women who have had treatment for breast or other cancers, including chemotherapy, aromatherapy blockers, and pelvic radiotherapy. Women who have had their ovaries removed surgically develop it early, because their oestrogen production stops abruptly.

\section*{Aetiology and Pathophysiology}
The vagina is lined by a layer of cells whose health depends on oestrogen. When oestrogen levels are high, the lining is thick, moist, and elastic, and it produces an acidic lubricating fluid. When oestrogen falls, the lining becomes thinner, paler, drier, and less flexible, and the natural lactobacilli that keep the vagina healthy become less numerous.

\begin{itemize}
    \item \textbf{Menopause and perimenopause:} the most common cause; oestrogen levels decline as the ovaries stop producing eggs
    \item \textbf{Breastfeeding:} high levels of the hormone prolactin suppress oestrogen, causing temporary dryness
    \item \textbf{Cancer treatments:} chemotherapy, pelvic radiotherapy, and hormonal treatments for breast cancer such as tamoxifen and aromatase inhibitors
    \item \textbf{Surgical removal of the ovaries:} causes an abrupt fall in oestrogen
    \item \textbf{Some contraceptives:} the progestogen-only pill, implant, and injection can reduce oestrogen effect and cause dryness in some women
    \item \textbf{Medicines:} some antidepressants, antihistamines, and cold remedies can cause generalised dryness, including the vagina
    \item \textbf{Other conditions:} Sjögren's syndrome, an autoimmune condition that dries the mucous membranes, and diabetes can contribute
    \item \textbf{Soaps and products:} harsh soaps, perfumed washes, and douching can irritate and dry the vulva and vagina
\end{itemize}

\section*{Clinical Features}
The symptoms of vaginal dryness are often understated by women but can be very distressing.

\begin{itemize}
    \item A feeling of dryness, tightness, or soreness in and around the vagina
    \item Itching, burning, or irritation that is not caused by infection
    \item Pain during sexual intercourse (dyspareunia), which can be severe
    \item Soreness or light bleeding after sex, because the thin lining is easily irritated
    \item Recurrent urinary symptoms, including passing urine more often or feeling a frequent urge
    \item Recurrent urinary tract infections or a burning sensation when urinating
    \item A generalised feeling that the genitals are uncomfortable or that something is not right
    \item Reduced interest in or avoidance of intimacy because of fear of pain
\end{itemize}

\section*{Differential Diagnosis}
\begin{itemize}
    \item Thrush (vaginal candidiasis), which also causes itching and soreness
    \item Bacterial vaginosis, which can cause discharge and irritation
    \item Contact dermatitis from soaps, fragrances, or laundry products
    \item Lichen sclerosus, a skin condition causing itching and scarring of the vulva
    \item A urinary tract infection, which causes similar urinary symptoms
    \item Vulvodynia, chronic vulval pain without a clear visible cause
    \item Sexually transmitted infections, which may cause vulval symptoms
    \item Diabetes, which can cause dryness and recurrent thrush
\end{itemize}

\section*{When to Seek Medical Care}
Women should seek medical advice if vaginal dryness causes pain during sex, discomfort in daily life, or emotional distress, or if it is accompanied by bleeding, discharge, or persistent itching, because those symptoms need to be checked to rule out infection or other causes. Women who have had breast cancer should seek specialist advice before using any hormonal treatment for dryness, because the choice of treatment depends on their individual situation.

\textbf{Contact your local emergency services for:}
\begin{itemize}
    \item Severe pain during sex or bleeding that is heavy or persistent
    \item A fever together with pelvic pain and discharge, which may indicate infection
    \item Any new symptom that feels alarming, such as sudden severe pelvic pain
\end{itemize}

\section*{Diagnosis and Investigations}
The diagnosis of vaginal dryness is usually made from the history and a simple examination.

\begin{itemize}
    \item \textbf{History:} symptoms, their duration, relationship to sex and the menstrual cycle, menopausal or breastfeeding status, medicines, cancer treatments, and previous gynaecological care
    \item \textbf{Physical examination:} gentle inspection of the vulva and vagina to assess thinning, paleness, and elasticity of the lining, and to exclude other causes such as infection or skin conditions
    \item \textbf{Swabs:} taken to exclude thrush or bacterial vaginosis when symptoms suggest infection
    \item \textbf{Hormone blood tests:} not usually needed for postmenopausal women, but may be useful when dryness occurs in younger women or with uncertain menopausal status
    \item \textbf{Discussion of contributing factors:} review of all medications and of the use of soaps and personal care products
\end{itemize}

\section*{Management}
Vaginal dryness responds very well to treatment, and the options are effective, safe, and widely available.

\begin{itemize}
    \item \textbf{Lubricants:} water-based or silicone-based lubricants used during sex, applied generously, to relieve friction and pain immediately
    \item \textbf{Vaginal moisturisers:} used regularly, two to three times a week, to keep the vaginal tissues hydrated; these are not just for use during sex
    \item \textbf{Vaginal oestrogen:} low-dose vaginal oestrogen as a cream, pessary, or ring, which is the most effective treatment; it is absorbed locally and is safe for most women, including many breast cancer survivors, though those women should discuss it with their oncology team
    \item \textbf{Systemic hormone replacement therapy (HRT):} can improve vaginal dryness as part of menopausal hormone therapy, when other menopausal symptoms are also present
    \item \textbf{Pelvic floor therapy and dilators:} helpful when dryness has led to tightening and muscle spasm around the vaginal opening, making penetration painful
    \item \textbf{Lifestyle measures:} gentle washing, avoiding irritants, and wearing comfortable cotton underwear
    \item \textbf{Regular sexual activity or stimulation:} keeping the tissues active helps maintain elasticity and moisture
\end{itemize}

\section*{Complications}
\begin{itemize}
    \item Painful intercourse that leads to avoidance of intimacy and relationship strain
    \item Reduced sexual desire and emotional distress, including anxiety about sex
    \item Recurrent urinary tract infections, because the thinning lining is more vulnerable to bacteria
    \item Bleeding after sex from the fragile lining
    \item Recurrent thrush, which is more common in an atrophic, less acidic vagina
    \item Impact on self-esteem and quality of life
    \item The risk, if untreated, of permanent tightening and narrowing of the vaginal entrance
\end{itemize}

\section*{Prognosis and Outcomes}
With treatment, vaginal dryness improves substantially in the great majority of women. Lubricants and moisturisers relieve symptoms for most, and vaginal oestrogen is highly effective, often restoring moisture and elasticity within weeks to a few months. For women who cannot use oestrogen, such as some breast cancer survivors, non-hormonal treatments and specialist guidance still provide significant relief. Dryness related to breastfeeding usually resolves when breastfeeding stops or as the hormones settle. The condition does not go away by itself in menopause without treatment, but it is very manageable, and women who seek help generally report a marked improvement in comfort, sexual enjoyment, and quality of life. There is no need to accept dryness as an unavoidable burden of ageing.

\section*{Prevention and Self-Care}
\begin{itemize}
    \item Use water-based or silicone-based lubricants generously during sex; silicone lubricants last longer
    \item Use a vaginal moisturiser two to three times a week even if you are not sexually active
    \item Avoid harsh soaps, perfumed washes, douches, and bubble baths in the genital area
    \item Wash with plain water or a mild, unscented cleanser
    \item Keep sexually active or stimulate the area regularly, as regular use helps maintain tissue health
    \item Wear cotton underwear and avoid very tight clothing
    \item If you have had breast cancer, discuss options with your oncology team before using any hormonal product
    \item See your doctor for a review if dryness is persistent or affects your quality of life
\end{itemize}

\section*{Sources and Further Reading}
The following organisations provide reliable information on vaginal dryness and menopausal vaginal health.

\begin{itemize}
    \item Jean Hailes for Women's Health. \textit{Vaginal dryness and menopause.} https://www.jeanhailes.org.au/
    \item Australasian Menopause Society. \textit{Vaginal symptoms and menopause.} https://www.menopause.org.au/
    \item NHS. \textit{Vaginal dryness and genitourinary syndrome of menopause.} https://www.nhs.uk/
    \item Mayo Clinic. \textit{Vaginal dryness: causes and treatment.} https://www.mayoclinic.org/
    \item MSD Manual. \textit{Atrophic vaginitis.} https://www.msdmanuals.com/
    \item MedlinePlus. \textit{Vaginal dryness.} https://medlineplus.gov/
\end{itemize}
